%% start of file `template.tex'.
%% Copyright 2006-2013 Xavier Danaux (xdanaux@gmail.com).
%
% This work may be distributed and/or modified under the
% conditions of the LaTeX Project Public License version 1.3c,
% available at http://www.latex-project.org/lppl/.

\documentclass[11pt,a4paper,sans]{moderncv}
\moderncvstyle{banking}
\moderncvcolor{blue}
\usepackage[utf8]{inputenc}
\usepackage[scale=0.75, margin=0.8in]{geometry}
\usepackage{import}
\usepackage{fontawesome}
\usepackage{footmisc}
\usepackage{dirtytalk}

\usepackage[sorting=ydnt]{biblatex}
\addbibresource{cite.bib}

\renewcommand*{\cventry}[7][.25em]{
  \begin{tabular*}{\maincolumnwidth}{l@{\extracolsep{\fill}}r}%
    % skip the company name and location line if both are empty, allowing one to show multiple roles within the same company without repeating it
    \ifboolexpr{%
      test {\ifstrempty{#4}}
      and
      test {\ifstrempty{#5}}}%
      {}%
      {{\bfseries #4} & {\bfseries #5}\\}%
    {\itshape #3\ifstrempty{#6}{}{\hspace{1cm} #6}} & {\itshape #2}\\%
  \end{tabular*}%
  \ifx&#7&%
  \else{\\%
    \begin{minipage}{\maincolumnwidth}%
      \small#7%
    \end{minipage}}\fi%
  \par\addvspace{#1}}

\name{\textcolor{color1}{A}\textcolor{gray}{lexandre}}{\textcolor{color1}{G}\textcolor{gray}{baguidi} \textcolor{color1}{A}\textcolor{gray}{ïsse}}
\title{CV}
\address{22 Rue de la République, 92150 Suresnes, France}{}{}
\phone[mobile]{+33 642426980}
\email{a.gbaguidiaisse@gmail.com}
% \homepage{agbaguidiaisse.com}
\social[linkedin]{agbaguidiaisse}
\social[github]{shumpaga}

\usepackage{background}
\usepackage{color}
\usetikzlibrary{calc}
\newsavebox\mybox
\sbox\mybox{%
\begin{tikzpicture}[remember picture]
    \fill[color1] (current page.north west) rectangle ($(current page.south west) +(1mm,0)$);
\end{tikzpicture}%
}
\backgroundsetup{
angle=0,
scale=1,
opacity=1,
color=black,
contents={%
\begin{tikzpicture}[remember picture, overlay]
    \node[anchor=west] at  (current page.west) {\usebox{\mybox}};
\end{tikzpicture}%
}
}

\usepackage{amsmath}
\newcommand{\CC }{\emph{ContextCapture}}
\newcommand{\CCMark }{\CC\footnotemark[1] }
\newcommand{\omg}{$\omega$-automata }
\usepackage{enumitem}
\definecolor{cusblue}{rgb}{0.0, 0.4, 0.6}

\begin{document}
\makecvtitle
\vspace{-1cm}
\begin{center}
  A la recherche d'un \textbf{sujet de thèse en informatique}.
\end{center}

\section{Profile}
J'ai choisi la recherche informatique comme spécialité durant mes études à l'EPITA.
Cela a joué un rôle clé dans mon travail chez Bentley Systems, m'enmenant à lire des papiers,
réimplémenter des algorithmes et déposer
\textbf{deux brevets d'invention}~\cite{patent.scanfinder1,patent.scanfinder2}.
Ayant compris que j'ai une forte passion pour la recherche, je souhaite faire un doctorat.
J'ai eu des expériences en:
\textbf{Vision par Ordinateur, Reconstruction 3D, Nuages de points, Optimization, Théorie des graphes,
Machines à Etats-Finis, Traitement de données, services web}.

\section{Experiences}
\footnotetext[1]{\CC\ est un logiciel de Reconstruction 3D à partir d'images ou de nuages de points décrivant une scène.}
\footnotetext[2]{Spot est une bibliothèque C++14 de manipulation de formules LTL, machines à états-finis et vérification formelle : \url{spot.lrde.epita.fr}.}
\begin{itemize}

\item{
  \cventry{Sep 2018 - auj.}{\textcolor{cusblue}{Ingénieur Logiciel}}{Bentley Systems -- Equipe "Reality Modeling"}{Paris, France}{\hspace{0.2cm}{\footnotesize Vision par ordinateur, Reconstruction 3D, Géométrie}}{}
  \vspace{-0.3cm}
  \cventry{Fev 2018 - Août 2018}{\textcolor{cusblue}{Stage ingénieur}}{}{}{\hspace{0cm}{\footnotesize Théorie des Graphes, Optimization, C++, Python, R, SQL}}{
    \vspace{0.1cm}
    Amélioration du logiciel \CCMark et son service cloud avec:
    \begin{itemize}[noitemsep,topsep=0.05cm,itemsep=0.05cm]
      \item Amélioration du support des nuages de point~\cite{techreport.contextcapture}:
        \begin{itemize}[noitemsep,topsep=0.05cm,itemsep=0.05cm]
          \item Un algorithme qui retrouve la position de tous les scanners d'un nuage de point multiple~\cite{patent.scanfinder1},
          \item Un algorithme qui détermine pour chaque point, le scanner l'ayant généré~\cite{patent.scanfinder2}.
          \item Différentes approches pour diminuer la sensibilité aux bruits des nuages de points.
        \end{itemize}
      \item Réécriture complète du méchanisme de jobs / tâches du logiciel sous forme de graphes de dépendances.
      \item Abstraction et convergence des des architectures (Desktop et Cloud).
    \end{itemize}
  }
}

\item{\cventry{Sep 2016 -- Jan 2018}{\textcolor{cusblue}{Etudiant chercheur}}{LRDE (Laboratoire de Recherche et Développement de l'EPITA)}{Paris, France}{{\footnotesize Machines à états-finis, Théorie des Graphes, Vérification formelle}}{
  \begin{itemize}[noitemsep,topsep=0.05cm,itemsep=0.05cm]
    \item Différentes contributions de code à Spot\footnotemark[2]:
      \begin{itemize}[noitemsep,topsep=0.05cm,itemsep=0.05cm]
        \item Amélioration d'un algorithme de minimisation de graphe en testant une approche incrémentale~\cite{techreport.minimization}
        \item Support des automates de co-B\"uchi (une classe de machines à états-finis avec des mots infinis)
        \item Test de persistence et de récurrence sur des formules LTL en utilisant les automates de co-B\"uchi
      \end{itemize}
    \item TD : Cours pratiques pour aider les étudiants à maîtriser les notions vues en cours 
  \end{itemize}
  }
}

\end{itemize}


\section{Education}
\begin{itemize}
\setlength\itemsep{-0.2cm}
\item{\cventry{Sep 2013 -- Sep 2018}{Ingénieur en informatique}{EPITA -- Spécialisation en Recherche Informatique}{Paris, France}{}{}}
\item{\cventry{Jan 2015 -- Mai 2015}{Cours de Systèmes temps-réel embarqués}{Staffordshire University -- Semestre à l'étranger}{Stafford, England} {}{}}
\end{itemize}

\section{Compétences}
\setlength{\tabcolsep}{3pt}
\renewcommand{\arraystretch}{1.2}
\begin{tabular}{ll}
\textbf{Linguistiques:} & Français (Natif), Russe (Natif), Anglais\\
\textbf{Informatique:} & C, C++, C\#, Python, R, Javascript, React, NodeJS, Scala, Java, PostgreSQL, SQLite.\\
\textbf{OS:} & Linux / UNIX, Mac Os X, Windows.\\
\textbf{Autres Intérêts:} & Education, Enregistrement et traitement Audio (Mix, master), voyages.
\end{tabular}

\printbibliography[type=patent,title={Brevets déposés}]
\printbibliography[type=report,title={Rapports écrits}]

\end{document}